% ============================================================
%  Chapter 6 -- Conclusion
% ============================================================
\chapter{Conclusion and Future Work}
\label{ch:conclusion}

% ============================================================
\section{Summary}

This thesis addressed the problem of geometric oversmoothing in learned
point cloud codecs, with a specific focus on PCGCv2 as a representative system.

We began by observing that standard aggregate quality metrics mask the
spatially non-uniform nature of codec artifacts.
To expose this, we introduced the \emph{curvature-stratified PSNR metric},
which partitions the original point cloud by local surface variation and
computes quality separately for edge (high-curvature) and flat (low-curvature)
regions.
Applied to PCGCv2 across the 8iVFB dataset, the metric revealed a consistent
and growing degradation gap of 2--4.5~dB, widening with rate --- a structural
bias of the learned codec that standard metrics do not detect.

With this diagnostic, we designed a post-processing network to correct the
identified distortions without modifying the codec.
The architecture --- a three-level sparse U-Net with a FiLM-conditioned
displacement head --- was shaped by three requirements derived from the
oversmoothing analysis:
(a) spatial awareness of surface geometry, provided by the U-Net's
multi-scale receptive field;
(b) rate adaptation, provided by the FiLM head;
(c) a training loss that avoids zero-inflation collapse, provided by the
dynamic Chamfer Distance.

The final model achieves an average edge PSNR improvement of
$\mathbf{+2.06}$~dB at the lowest evaluated rate,
$\mathbf{+0.58}$~dB at medium rate, and
$\mathbf{+0.27}$~dB at high rate,
across all four 8iVFB sequences and 400 frames.
These gains are consistent across all four sequences, including three not
seen during training.

Zero-shot transfer experiments demonstrate that the model extracts
corrections general enough to improve an unseen codec (Unicorn) at extreme
compression rates, where oversmoothing is severe enough to be codec-agnostic.

% ============================================================
\section{Contributions}

To summarise, the concrete contributions of this thesis are:

\begin{enumerate}
  \item \textbf{Curvature-stratified PSNR.}
    A diagnostic metric that separates codec evaluation by surface geometry,
    enabling targeted assessment of edge-region quality.
    The implementation is released as part of the evaluation pipeline.

  \item \textbf{Oversmoothing characterisation.}
    A systematic quantitative study showing that PCGCv2's oversmoothing gap
    grows with rate, is sequence-dependent, and is absent in G-PCC, providing
    evidence that it is a structural property of the learned encoder.

  \item \textbf{Dynamic Chamfer Distance loss.}
    A per-patch matching loss that avoids both zero-inflation collapse and
    nearest-neighbour noise, enabling successful training of a displacement
    regressor where standard losses fail.

  \item \textbf{FiLM-conditioned multi-rate post-processor.}
    A single lightweight model that serves multiple rate-distortion operating
    points, achieving better performance than per-rate models through
    beneficial multi-rate regularisation.

  \item \textbf{Evidence for cross-codec oversmoothing.}
    Zero-shot transfer results showing that at extreme compression, geometric
    oversmoothing is a codec-agnostic failure mode that a learned corrector
    can address without retraining.
\end{enumerate}

% ============================================================
\section{Limitations}

\paragraph{Single-codec training.}
The model is trained exclusively on PCGCv2-decoded clouds.
At moderate rates for other codecs (e.g., Unicorn R2--R4), transfer is
negative.
A codec-agnostic post-processor would require either multi-codec training
or a codec identification mechanism.

\paragraph{Displacement-only formulation.}
The network can only shift existing points, not add or remove them.
At very low bitrates, the decoded cloud may be missing entire surface regions
(point dropping), which displacement corrections cannot recover.

\paragraph{Dataset scope.}
All training and primary evaluation is on the 8iVFB dataset: dense,
full-body human captures at $G=1024$.
Performance on sparse outdoor LiDAR scans, mechanical parts, or lower-resolution
captures is not evaluated.

\paragraph{Computational cost.}
Post-processing adds an inference step with latency proportional to point
cloud size.
For streaming applications, this latency must be bounded; the current
implementation is not optimised for real-time use.

% ============================================================
\section{Future Work}
\label{sec:future-work}

The results open several directions for follow-up.
We group them roughly by how much new work each would require.

\subsection*{Near-Term Extensions}

\paragraph{Curriculum fine-tuning.}
The current model trains all rates jointly from the start.
A curriculum approach --- pre-training the backbone on r2 alone, then
fine-tuning the FiLM head on r4 and r6 --- may further improve mid-to-high
rate performance.
The intuition: the r2 displacement signal is the richest, so the backbone
develops strong edge-detection features from r2 data.
Fine-tuning the FiLM head separately for each rate, with the backbone frozen,
could allow the head to specialise without interference from competing
rate objectives.
Preliminary experiments with this curriculum order showed promising
convergence on validation loss at r4, but the approach was not fully evaluated
for the final paper.

\paragraph{Inference-time rate adaptation.}
The current model is conditioned on the actual BPP of each decoded cloud,
which requires metadata from the codec.
In deployment settings where the BPP is unknown, the model could estimate it
from the decoded cloud statistics (e.g., decoded point count divided by the
number of transmitted bytes, reconstructed from the occupancy pattern).
A blind BPP estimator would make the post-processor completely black-box,
requiring only the decoded point cloud.

\paragraph{Faster inference via learned patch routing.}
The sliding-window patch inference processes all patches uniformly.
At high bitrates, most patches are already close to the original surface and
require only near-identity corrections.
A lightweight ``patch quality'' predictor (e.g., based on local point density
variance) could route patches to skip inference when the estimated correction
is below a threshold, reducing inference time at high bitrates by
$50$--$70\%$ with negligible quality loss.

\subsection*{Medium-Term Research Questions}

\paragraph{Point addition via occupancy prediction.}
The displacement formulation can only shift existing points, not add new ones.
At very low bitrates (below r1), PCGCv2 drops entire surface patches, leaving
regions empty that no displacement can recover.
Extending the displacement head with an additional occupancy prediction branch
--- predicting whether new points should be inserted in empty voxels near the
current surface --- would address this limitation.
This requires a different loss: the Chamfer Distance extends naturally to
handle variable-size output sets (the CD itself does not require
fixed cardinality), but the occupancy branch introduces a combinatorial
structure that needs careful regularisation to avoid spurious insertions.

\paragraph{Attribute (colour) post-processing.}
This thesis focuses entirely on geometry.
Learned codecs also oversmooth colour attributes at high-curvature regions:
colour boundaries (e.g., the edge between a face and hair) are blurred in
the decoded cloud by the same occupancy-based encoding mechanism that blurs
geometry.
A joint geometry-and-colour post-processor using the geometric features as
context for colour refinement is a natural extension.
The shared backbone would extract edge-aware features; separate heads would
predict geometric displacements and colour residuals.
The loss would combine geometric Chamfer Distance and colour mean squared
error, with an optional edge-aware perceptual weight.

\paragraph{Temporal coherence in video sequences.}
The current post-processor is applied frame-by-frame, independently.
For volumetric video sequences at 10 fps, adjacent frames share approximately
90\% of surface geometry (a human subject moves slowly).
A temporally coherent post-processor could condition on the refined cloud
from the previous frame, propagating high-quality corrections forward in time
and reducing temporal flickering artifacts that are perceptually disturbing
in volumetric video playback.
The FiLM conditioning is a natural extension point: the temporal offset (frame
number within the GOP) could be appended to the rate scalar.

\paragraph{Multi-codec training.}
At intermediate bitrates (Unicorn R2--R4), the current model applied to
Unicorn-decoded clouds produces negative PSNR changes.
This is because Unicorn's distortion patterns differ from PCGCv2's at
non-extreme rates.
Training on a mixture of PCGCv2- and Unicorn-decoded clouds, with a codec
embedding (e.g., a one-hot or learned codec ID) as an additional conditioning
signal, would extend generalisation to intermediate rates for multiple codecs.
The FiLM head already accepts a vector conditioning input; appending the
codec embedding to the rate scalar requires minimal architectural changes.

\subsection*{Longer-Term Directions}

\paragraph{Evaluation on LiDAR sequences.}
All training and evaluation in this thesis uses the 8iVFB dataset: dense,
full-body human captures at $G = 1024$ with approximately $10^6$ points per
frame.
LiDAR point clouds (e.g., SemanticKITTI, nuScenes) have quite different
geometric properties: sparse, outdoor, with large empty regions
and a distribution of object types (cars, buildings, pedestrians).
Whether the curvature-stratified metric and the displacement post-processor
transfer to this domain is an open question.
If they do not, a LiDAR-specific training set and possibly a modified
curvature threshold (LiDAR surfaces have lower curvature density) would be
needed.

\paragraph{Perceptual quality metric integration.}
The curvature-stratified PSNR introduced in this thesis is a step toward
perceptually motivated evaluation.
A fully perceptual metric would require psychophysical experiments with
human observers, as standard in image quality assessment (SSIM, LPIPS).
For volumetric video, where the viewer can rotate the viewpoint, the
perceptual weighting of surface regions depends on the viewing direction ---
a direction-dependent metric that accounts for self-occlusion and viewing
probability would be more predictive of subjective quality.
Such a metric could then be used as the training objective for a
perception-driven post-processor.

\paragraph{Comparison with newer codecs and MPEG standardisation.}
SparsePCGC~\cite{wang2023sparsepcgc} and Unicorn~\cite{wang2025unicorn}
are the current state of the art.
Future developments in MPEG Immersive Coding (e.g., Video-based Point Cloud
Compression v2) may exhibit different oversmoothing profiles.
Applying the curvature-stratified evaluation framework systematically to all
these systems would establish whether geometric oversmoothing is a property
of the occupancy-grid paradigm in general, or specific to PCGCv2's
coarse-to-fine architecture.
This would contextualise the PCGCv2 findings within the broader trajectory
of the field and inform which architectural choices are responsible for the
edge-region quality gap.

\paragraph{Neural radiance field (NeRF) and 3D Gaussian codecs.}
Recent neural scene representations (NeRF, 3D Gaussian Splatting) are
beginning to be used for volumetric video compression.
These representations do not use explicit point clouds and are not directly
susceptible to geometric oversmoothing in the same way.
However, they exhibit their own characteristic artifacts at low bitrates
(blurriness, floaters, missing thin structures).
Developing analogous diagnostic metrics for these representations, and
learned post-processors that correct their specific failure modes, extends
the methodology developed here to a broader class of codecs.
